\chapter{La match analysis}
\label{cap:matchanalysis}

\textbf{Match analysis}: la branca della scienza dello sport che studia che cosa accade durante la
partita. Serve al preparatore per una ragione precisa: \textbf{dobbiamo sapere che cosa valutare},
e questo ce lo indica la match analysis. È il ponte fra i principi di tecnica e tattica del
cap.~\ref{cap:tecnicatattica} e la scelta dei test di valutazione funzionale.

Tutti gli studi presentati sono pubblicati su \textbf{riviste internazionali impattate}: la scelta
è deliberata, perché accanto all'esperienza di campo servono \textbf{riferimenti certi}.

% ---------------------------------------------------------------------------
\section{Che cosa studia la match analysis}

L'impianto di riferimento è di \textbf{D'Ottavio S.} La numerazione della slide passa direttamente
dal punto 1 al punto 4: i punti 2 e 3 non compaiono nel materiale.

\rubrica{1. Studio della performance in gara}
\begin{enumerate}
  \item \textbf{distanza totale percorsa} e zone a maggiore densità di spostamento (ripartizione in
    zone);
  \item differenze tra \textbf{primo e secondo tempo};
  \item distanza ai vari \textbf{range di velocità}: quanto tempo il giocatore trascorre in un dato
    range, o quanta distanza percorre in quel range;
  \item profilo delle attività a \textbf{intervalli di tempo} (per esempio ogni 15 minuti), così da
    confrontare l'inizio e la fine di un tempo di gioco, o i soli ultimi minuti;
  \item \textbf{profilo fisiologico}: FC, LA, FFA, ammonio, ormoni;
  \item profilo delle attività svolte ad \textbf{alta intensità con e senza palla};
  \item numero di attività di \textbf{sprint}, accelerazioni, accelerazioni sprint, decelerazioni,
    e la loro intensità;
  \item analisi dei \textbf{cambi di direzione} e valutazione degli angoli utilizzati;
  \item studio delle influenze dei diversi \textbf{sistemi di gioco}, cioè di come le indicazioni
    tattiche dell'allenatore influiscano sul profilo cinematico e fisiologico;
  \item analisi delle differenze relative al \textbf{ruolo} in campo.
\end{enumerate}
Conoscere questi valori permette di \textbf{selezionare i test di valutazione} più adatti e
sport-specifici, e di rendere l'allenamento in campo il più funzionale possibile alla prestazione.

\rubrica{4. Studi comparativi}
\begin{itemize}
  \item \textit{top level} contro \textit{low level};
  \item calcio maschile e calcio femminile;
  \item giovani e adulti;
  \item giovani \textbf{elite} e giovani \textbf{sub elite}.
\end{itemize}

\rubrica{5. Studio delle relazioni ecologiche fra test e performance in gara}
Una volta nota la performance di gara, si verifica se il test proposto abbia \textbf{relazioni
strette} con quello che succede davvero in partita. Mette in relazione i \textbf{parametri
cinematici e fisiologici rilevati in gara} con le \textbf{qualità fisiche rilevate con i test} da
campo e da laboratorio --- per esempio \textit{high intensity activity}, RSA, Yo-Yo test, CMJ.

\rubrica{6. Organizzazione di attività allenanti specifiche dettate dagli studi precedenti}
L'obiettivo è \textbf{riprodurre in allenamento le intensità fisiche e fisiologiche della gara}.
\begin{itemize}
  \item \textbf{percorsi misti}: percorsi in cui si alternano tipologie di corsa a intensità
    diverse, con l'inserimento di cambi di direzione o della \textbf{palla};
  \item \textbf{small sided games}: partite a campo ridotto. Nel calcio a 5 si definiscono tali
    tutti i giochi \textbf{fino al 4 contro 4}, cioè con un numero di giocatori inferiore a quello
    della partita. Sono molto usati perché riproducono la prestazione di gara \textbf{dentro una
    situazione di gara}: c'è il pallone, c'è l'avversario, e quindi c'è un'\textbf{imprevedibilità
    degli eventi} molto simile a quella che il giocatore dovrà risolvere in partita. I parametri
    che ne orientano l'intensità:
    \begin{enumerate}
      \item \textbf{spazio}: lo stesso 4 contro 4 su 30 $\times$ 20\,m o su 20 $\times$ 10\,m dà
        risposte cinematiche e fisiologiche diverse;
      \item \textbf{numero di giocatori};
      \item \textbf{regole} tecniche e tattiche $\rightarrow$ intensità;
      \item \textbf{tempo} (durata);
      \item \textbf{rapporto tempo attivo / tempo di recupero};
    \end{enumerate}
  \item \textbf{training di corsa e andature tecniche} dirette all'attivazione di energia aerobica,
    anaerobica e combinata neuromuscolare;
  \item \textbf{training prettamente neuromuscolare}.
\end{itemize}

% ---------------------------------------------------------------------------
\section{Perché nel futsal si usa il video e non il GPS}

È il vincolo che condiziona tutta la letteratura di questo sport.

\rubrica{Il GPS non è utilizzabile}
Il calcio a 5 è uno sport \textbf{indoor}: la struttura del palazzetto \textbf{impedisce o disturba
il segnale GPS}, che quindi non può essere impiegato. Negli sport outdoor invece il GPS è la
tecnologia corrente: applicato agli sport di squadra a partire dal \textbf{2000} con frequenze di
campionamento bassissime, oggi fornisce un dato molto più accurato, e la maggior parte degli studi
sul profilo cinematico lo utilizza.

\rubrica{Il video tracking e il suo costo}
Resta come unico metodo possibile la \textbf{videoanalisi}:
\begin{itemize}
  \item richiede \textbf{almeno tre telecamere} che agiscono contemporaneamente, poste nei punti
    più alti rispetto al campo;
  \item le riprese vengono \textbf{triangolate} per derivare le traiettorie dei giocatori;
  \item i costi sono \textbf{molto elevati}, e rilevare questi dati è quindi estremamente
    difficile.
\end{itemize}

\rubrica{I due tipi di software}
\begin{itemize}
  \item \textbf{automatici}: i più avanzati e i più costosi; svolgono il lavoro da soli;
  \item \textbf{semi-automatici}: è l'\textbf{operatore} a tracciare su un reticolo le traiettorie
    del giocatore, e il software restituisce i dati cinematici --- distanza percorsa, velocità
    raggiunta, intensità dell'accelerazione.
\end{itemize}

\rubrica{La conseguenza}
La difficoltà tecnica spiega perché fino al \textbf{2013--2014} fossero pochissimi gli studi che
avessero analizzato la prestazione in questo sport --- molti di più quelli sull'\textbf{epidemiologia
degli infortuni}.

% ---------------------------------------------------------------------------
\section{I range di velocità}

La ripartizione della distanza percorsa in fasce di velocità è dovuta a \textbf{Castagna e coll.,
2009} ed è stata assunta come \textbf{standard} da tutta la letteratura successiva: la si ritrova
identica in Makaje (2012), Vieira (2016) e Beato (2017).

\rubrica{La versione a cinque bande}
\begin{itemize}
  \item $V_1 \leq 6{,}0$\,km/h --- \textit{standing and walking};
  \item $6{,}1 < V_2 \leq 12{,}0$\,km/h --- \textit{low-intensity running};
  \item $12{,}1 < V_3 \leq 15{,}4$\,km/h --- \textit{medium-intensity running};
  \item $15{,}5 < V_4 \leq 18{,}3$\,km/h --- \textit{high-intensity running};
  \item $V_5 > 18{,}4$\,km/h --- \textit{sprinting}.
\end{itemize}

\rubrica{La versione a sei bande}
Nello studio di Castagna stesso la banda più bassa è \textbf{divisa in due} --- ed è proprio quella
che nella versione a cinque porta l'etichetta \textit{standing and walking}:
\begin{itemize}
  \item \textit{standing}: 0--0,4\,km/h, giocatore fermo;
  \item \textit{walking}: 0,5--6\,km/h;
  \item \textit{low-intensity running}: 6,1--12\,km/h;
  \item \textit{medium-intensity running}: 12,1--15,4\,km/h;
  \item \textit{high-intensity running}: $> 15{,}5$\,km/h;
  \item \textit{sprinting}: $> 18{,}3$\,km/h.
\end{itemize}

\rubrica{Il limite dell'analisi per velocità: serve l'accelerazione}
È una riflessione da fare \textbf{prima} di leggere qualunque profilo di attività.
\begin{itemize}
  \item \textbf{l'obiezione}: il campo di 40 $\times$ 20\,m \textbf{limita di fatto} la possibilità
    del giocatore di raggiungere velocità elevate;
  \item \textbf{la risposta}: l'analisi per range di velocità resta \textbf{assolutamente corretta}
    e ci fa capire quale sia la prestazione del giocatore di calcio a 5, specialmente nei momenti di
    massimo sforzo;
  \item \textbf{ma non basta}: per un'analisi ancora più corretta serve il dato
    dell'\textbf{accelerazione}, cioè \textit{quanto velocemente} il giocatore ha raggiunto quella
    velocità;
  \item \textbf{l'esempio}: uno sprint massimale da fermo, ma interrotto a 10\,m per la presenza
    della palla, dell'avversario o per il limite del campo, ha un'\textbf{accelerazione molto
    forte} ed è \textbf{costato molto dal punto di vista energetico}, pur senza aver raggiunto
    velocità elevate. Letto con il solo criterio della velocità, quello sprint sparisce.
\end{itemize}

% ---------------------------------------------------------------------------
\section{Il profilo della gara}

\subsection{Bueno e coll., 2014 --- la distanza percorsa}

\textit{Analysis of the distance covered by Brazilian professional futsal players during official
matches}, Sports Biomechanics, 13(3), 230--240.

\rubrica{Scopo, metodo e protocollo}
Misurare la \textbf{distanza percorsa} con video tracking a tre telecamere. \textbf{5 partite
ufficiali}, $n = 93$ giocatori; range di velocità nella versione a cinque bande.

\rubrica{Risultati --- le distanze}
L'analisi distingue i momenti in cui la \textbf{palla è in gioco} da quelli in cui non lo è: nel
calcio a 5 il cronometro si ferma quando la palla esce, e questa distinzione è il contributo
principale dello studio.
\begin{itemize}
  \item \textbf{distanza totale}: 3133,2\,m, con lieve differenza fra i due tempi (primo 1710,6\,m,
    secondo 1635,9\,m);
  \item \textbf{distanza \textit{in play}}: 2133,9\,m --- circa \textbf{un chilometro in meno} di
    quanto direbbe l'analisi tradizionale;
  \item \textbf{distanza \textit{out of play}}: 1028,5\,m.
\end{itemize}

\rubrica{Risultati --- la distanza relativa (m/min)}
È qui che la distinzione pesa di più:
\begin{itemize}
  \item \textit{in play}: $\approx 136$\,m/min nel primo tempo, $\approx 129$\,m/min nel secondo;
  \item \textbf{partita intera}: $\approx 97$--98\,m/min nel primo tempo, $\approx 90$\,m/min nel
    secondo.
\end{itemize}
Analizzando la sola partita nel suo complesso si otterrebbe un dato \textbf{nettamente inferiore} a
quello reale, con conseguenze dirette sulla scelta dei test e sul \textbf{volume di allenamento} a
cui sottoporre i giocatori.

\rubrica{Risultati --- la ripartizione percentuale}
Percentuale della distanza percorsa in ciascuna fascia di velocità. Valori in \textbf{mediana
(IQR)}; l'asterisco indica differenza significativa rispetto al primo tempo ($p < 0{,}05$).

\begin{center}
\footnotesize
\begin{tabular}{|l|l|c|c|c|}
\hline
\textbf{Condizione} & \textbf{Fascia di velocità} & \textbf{1\textsuperscript{o} tempo} & \textbf{2\textsuperscript{o} tempo} & \textbf{$p$} \\
\hline
\multirow{5}{*}{\textit{In play}}
  & \textit{standing and walking}      & 16,2 (5,7)  & 19,3 (8,3)*  & $<0{,}01$ \\
  & \textit{low-intensity running}     & 41,9 (5,3)  & 42,1 (5,4)   & 0,69 \\
  & \textit{medium-intensity running}  & 20,1 (4,2)  & 17,8 (5,1)*  & $<0{,}01$ \\
  & \textit{high-intensity running}    & 10,3 (3,5)  & 9,6 (3,4)*   & $<0{,}01$ \\
  & \textit{sprinting}                 & 10,1 (6,1)  & 9,9 (5,0)    & 0,49 \\
\hline
\multirow{5}{*}{\textit{Out of play}}
  & \textit{standing and walking}      & 52,4 (11,9) & 55,4 (15,2)  & 0,72 \\
  & \textit{low-intensity running}     & 33,1 (8,0)  & 32,9 (11,1)  & 0,44 \\
  & \textit{medium-intensity running}  & 8,1 (5,9)   & 8,7 (5,5)    & 0,55 \\
  & \textit{high-intensity running}    & 2,1 (2,4)   & 3,1 (3,2)*   & $<0{,}01$ \\
  & \textit{sprinting}                 & 1,5 (2,8)   & 1,7 (3,0)    & 0,29 \\
\hline
\multirow{5}{*}{\textbf{Partita intera}}
  & \textit{standing and walking}      & 28,0 (6,1)  & 30,8 (6,7)*  & $<0{,}01$ \\
  & \textit{low-intensity running}     & 39,0 (5,0)  & 38,7 (4,0)   & 0,92 \\
  & \textit{medium-intensity running}  & 16,4 (3,4)  & 15,4 (3,4)*  & $<0{,}01$ \\
  & \textit{high-intensity running}    & 8,0 (2,4)   & 7,5 (2,0)*   & $<0{,}01$ \\
  & \textit{sprinting}                 & 7,6 (4,3)   & 7,2 (2,7)    & 0,32 \\
\hline
\end{tabular}
\end{center}

\begin{itemize}
  \item sommando alta intensità e sprint, l'\textit{in play} vale circa il \textbf{20\%}
    dell'attività, contro il \textbf{16\%} circa che si otterrebbe considerando la partita intera:
    una differenza percentuale notevole;
  \item nel secondo tempo, sia \textit{in play} sia sulla partita intera, si osserva un
    \textbf{processo di affaticamento}: aumenta la distanza percorsa camminando o correndo a bassa
    intensità e diminuisce, seppur di poco, l'attività ad alta intensità.
\end{itemize}

\rubrica{Gli istogrammi di velocità}
Mostrano la stessa cosa in forma grafica: in grigio l'attività \textit{out of play}, in nero quella
\textit{in play}.
\begin{itemize}
  \item colorando tutti gli istogrammi \textbf{dello stesso colore} si concluderebbe che la maggior
    parte dell'attività si svolge da fermo o camminando a bassa velocità;
  \item guardando i soli istogrammi \textbf{neri} si vede invece che la prestazione si svolge
    pochissimo a bassa velocità, e per la maggior parte \textbf{dagli 8\,km/h in su}.
\end{itemize}

\rubrica{Take home message}
\begin{itemize}
  \item la distanza percorsa \textbf{decrementa} tra primo e secondo tempo;
  \item nel secondo tempo \textbf{aumenta la distanza percorsa camminando};
  \item importanza dell'analisi \textbf{\textit{in play} -- \textit{out of play}} negli studi di
    match analysis nel futsal.
\end{itemize}

\subsection{Castagna e coll., 2009 --- le richieste della gara}

\textit{Match demands of professional Futsal: A case study}, Journal of Science and Medicine in
Sport, 12, 490--494.

\rubrica{Perché è uno studio di riferimento}
È uno dei \textbf{primi} ad aver analizzato la performance di gara nel calcio a 5, in un'epoca in
cui gli studi su questo sport erano pochissimi. Viene citato costantemente, e la sua ripartizione
in range di velocità è diventata lo \textbf{standard} della letteratura successiva.

\rubrica{Scopo e metodo: carico esterno e carico interno}
Analisi della prestazione \textbf{cinematica e fisiologica} insieme. Tradotto nel gergo
dell'allenamento:
\begin{itemize}
  \item con la \textbf{videoanalisi} gli autori studiano il \textbf{carico esterno}, cioè quello a
    cui i giocatori sono sottoposti durante la partita;
  \item con \textbf{FC, consumo di ossigeno e lattato ematico} monitorano il \textbf{carico
    interno}, cioè tutte le risposte fisiologiche che i giocatori mettono in atto per sostenere
    quell'attività.
\end{itemize}

\rubrica{Risultati --- cinematici}
\begin{itemize}
  \item \textbf{distanza percorsa}: 121\,m/min;
  \item \textbf{distanza media dello sprint}: 10,5\,m (6,2--14,8). È fortemente influenzata dalla
    presenza dell'avversario, dalla palla e soprattutto dal \textbf{limite delle dimensioni del
    campo};
  \item \textbf{frequenza degli sprint}: sequenze di \textbf{3--4 sprint} consecutivi ogni
    \textbf{20--30 secondi};
  \item l'attività ad alta intensità si attesta fra il \textbf{15 e il 20\%} della prestazione
    totale, in accordo con lo studio di Bueno.
\end{itemize}
Il dato sulle sequenze orienta sia l'allenamento sia la valutazione: il giocatore altamente
performante è quello che riesce a svolgere queste sequenze \textbf{recuperando nel minor tempo
possibile}, e saper valutare questa capacità è essenziale.

\rubrica{Risultati --- fisiologici}
\begin{itemize}
  \item $VO_2$: \textbf{76\% del $VO_{2}max$} (59--92), con range ampio;
  \item HR: \textbf{90\% della $HR_{max}$} (84--96). Qui il range si assottiglia, ed è quindi un
    valore che si può prendere come \textbf{pienamente valido};
  \item il tempo trascorso \textbf{sopra l'85\%} --- e anche sopra l'80\% --- della frequenza
    cardiaca massima è estremamente rilevante;
  \item lattato ematico: \textbf{5,3\,mmol/l} (1,1--10,4), con range molto ampio, probabilmente per
    la disparità di tempo di gioco fra i singoli giocatori.
\end{itemize}

\rubrica{Perché il lattato non è alto}
5,3\,mmol/l \textbf{non è un valore elevato}, e questo \textbf{non significa che non venga prodotto
acido lattico}. Trattandosi di uno sport \textbf{intermittente}, in cui alle fasi di alta intensità
corrispondono fasi di recupero --- da fermi, camminando o correndo a bassa intensità --- il lattato
prodotto viene \textbf{riutilizzato durante i recuperi}: non ci si trova quasi mai in presenza di un
vero e proprio \textbf{accumulo}.

\rubrica{Take home message}
\begin{itemize}
  \item il futsal è un \textbf{regime intermittente ad alta intensità};
  \item il \textbf{meccanismo aerobico/anaerobico} è fortemente sollecitato;
  \item la \textbf{repeated sprint ability} --- la capacità di ripetere sprint e di recuperare
    velocemente per essere pronti a una nuova azione ad alta intensità --- è una capacità specifica
    e fondamentale del futsal;
  \item ne consegue la necessità di \textbf{sviluppare la potenza aerobica}.
\end{itemize}

% ---------------------------------------------------------------------------
\section{Gli sprint e le sequenze di sprint ripetuti}

\subsection{Caetano e coll., 2015}

\textit{Characterization of the Sprint and Repeated-Sprint Sequences Performed by Professional
Futsal Players, According to Playing Position, During Official Matches}, Journal of Applied
Biomechanics, 31, 423--429.

\rubrica{Scopo, metodo e protocollo}
Analisi delle \textbf{repeated sprint sequencies} con video tracking, distinte per \textbf{ruolo}.
\textbf{5 partite ufficiali}, $n = 97$ giocatori.
\begin{itemize}
  \item \textbf{sprint}: ogni evento sopra i $5{,}08$\,m/s, cioè circa \textbf{18--19\,km/h}, in
    linea con la soglia di \textit{sprinting} degli studi precedenti;
  \item \textbf{repeated sprint sequency}: \textbf{2 o più sprint consecutivi} con intervallo fra
    loro di 15\,s (RS15), 30\,s (RS30), 45\,s (RS45) o 60\,s (RS60).
\end{itemize}

\rubrica{Risultati --- le variabili dello sprint}
\textbf{Non ci sono differenze fra i ruoli} (difensore, laterale, pivot): si riportano quindi i
valori complessivi, in media (deviazione standard).

\begin{center}
\footnotesize
\begin{tabular}{|l|c|c|}
\hline
\textbf{Variabile} & \textbf{1\textsuperscript{o} tempo} & \textbf{2\textsuperscript{o} tempo} \\
\hline
Distanza percorsa per sprint (m)      & 13,3 (5,7)  & 14,0 (6,5) \\
Durata (s)                            & 3,1 (1,2)   & 3,2 (1,3)* \\
Velocità di picco (m/s)               & 5,9 (0,7)   & 5,9 (0,7) \\
Velocità iniziale (m/s)               & 1,4 (1,2)   & 1,4 (1,2) \\
Tempo di recupero fra sprint (s)      & 55,3 (60,5) & 63,2 (71,6) \\
Sprint al minuto                      & 0,9 (0,4)   & 0,8 (0,4) \\
\hline
\end{tabular}
\end{center}
\footnotesize *\,unica differenza significativa rispetto al primo tempo ($p < 0{,}05$): la
\textbf{durata dello sprint}. Laterali e pivot raggiungono nel secondo tempo un picco di 6,0
(0,8)\,m/s.
\normalsize

\begin{itemize}
  \item la \textbf{distanza per sprint} è di circa 13\,m, contro i circa 10 di Castagna;
  \item il \textbf{picco di velocità} si ferma intorno ai 6\,m/s per le limitazioni di spazio, e
    perché generalmente non si parte da fermi --- come conferma la \textbf{velocità iniziale} di
    1,4\,m/s;
  \item si effettua circa \textbf{uno sprint al minuto}.
\end{itemize}

\rubrica{Risultati --- la frequenza delle sequenze}
Gli autori hanno \textbf{accorpato le cinque partite} invece di ricavare un valore medio per
partita.

\begin{center}
\footnotesize
\begin{tabular}{|l|c|c|c|c|}
\hline
\textbf{Sequenza} & \textbf{2 sprint} & \textbf{3 sprint} & \textbf{4 sprint} & \textbf{Totale (\%)} \\
\hline
RS15        & 368 (32,7\%) & 74 (6,6\%)  & 17 (1,5\%) & 40,8 \\
RS30        & 276 (24,6\%) & 22 (2,0\%)  & 8 (0,7\%)  & 27,2 \\
RS45        & 167 (14,9\%) & 22 (2,0\%)  & 3 (0,3\%)  & 17,1 \\
RS60        & 155 (13,8\%) & 12 (1,1\%)  & 0 (0,0\%)  & 14,9 \\
\hline
\textbf{Totale (\%)} & \textbf{85,9} & \textbf{11,6} & \textbf{2,5} & \textbf{100,0} \\
\hline
\end{tabular}
\end{center}

\begin{itemize}
  \item le sequenze di \textbf{due soli sprint} coprono l'\textbf{85,9\%} dei casi, e la
    combinazione più frequente in assoluto è \textbf{2 sprint con 15 secondi di recupero}
    (32,7\%);
  \item \textbf{più aumenta il tempo di recupero, più diminuisce la frequenza}: è la
    testimonianza dell'intensità dell'esercizio;
  \item le sequenze di tre e soprattutto di quattro sprint sono rare (11,6\% e 2,5\%), e con RS60
    non si presentano mai.
\end{itemize}

\rubrica{Take home message}
\begin{itemize}
  \item \textbf{non ci sono differenze fra i ruoli};
  \item le sequenze più frequenti sono di \textbf{2 sprint con 15 secondi di intervallo};
  \item \textbf{distanza media} 13,3\,m: dato prezioso per l'allenamento dell'accelerazione, della
    decelerazione e della velocità;
  \item la \textbf{durata dello sprint} subisce un lieve incremento fra primo e secondo tempo ---
    differenza piccola, ma che può indirizzare allenamento e valutazione.
\end{itemize}

% ---------------------------------------------------------------------------
\section{L'influenza del livello di qualificazione}

\subsection{Makaje e coll., 2012}

\textit{Physiological demands and activity profiles during futsal match play according to
competitive level}, Journal of Sports Medicine and Physical Fitness, 52, 366--374.

\rubrica{Scopo, metodo e protocollo}
Stabilire \textbf{richiesta fisiologica} e \textbf{profilo delle attività} confrontando
\textbf{elite} e \textbf{non elite}, distinguendo inoltre i \textbf{giocatori di movimento} dai
\textbf{portieri}. Video tracking per i parametri cinematici, FC e lattato ematico per quelli
fisiologici; range di velocità identici a quelli di Castagna e D'Ottavio.

Giocatori di movimento: $n = 12$ elite e $n = 12$ amatori. Portieri: $n = 3$ e $n = 3$. Valori in
\textbf{media $\pm$ deviazione standard}; l'asterisco indica differenza significativa fra i gruppi
($p < 0{,}05$).

\rubrica{Risultati --- tempo trascorso nei range di velocità (\%)}

\begin{center}
\footnotesize
\begin{tabular}{|l|c|c|c|c|}
\hline
 & \multicolumn{2}{c|}{\textbf{Giocatori di movimento}} & \multicolumn{2}{c|}{\textbf{Portieri}} \\
\hline
\textbf{Tempo trascorso} & \textbf{Elite} & \textbf{Amatori} & \textbf{Elite} & \textbf{Amatori} \\
\hline
\textit{Standing}              & 4,2 $\pm$ 1,1  & 6,9 $\pm$ 1,7*  & 8,2 $\pm$ 2,7  & 12,4 $\pm$ 2,3* \\
\textit{Walking}               & 26,1 $\pm$ 1,8 & 27,8 $\pm$ 2,2* & 51,0 $\pm$ 5,2 & 50,6 $\pm$ 4,2 \\
\textit{Jogging}               & 18,0 $\pm$ 1,6 & 17,2 $\pm$ 1,3* & 15,2 $\pm$ 4,2 & 14,1 $\pm$ 4,3 \\
\textit{Low-speed running}     & 19,4 $\pm$ 2,4 & 18,6 $\pm$ 1,7* & 10,4 $\pm$ 2,2 & 9,1 $\pm$ 2,5 \\
\textit{Moderate-speed running}& 17,1 $\pm$ 2,2 & 16,2 $\pm$ 2,6* & 8,3 $\pm$ 1,5  & 7,8 $\pm$ 1,5 \\
\textit{High-speed running}    & 8,7 $\pm$ 1,3  & 7,7 $\pm$ 1,6*  & 4,7 $\pm$ 1,3  & 4,2 $\pm$ 1,3 \\
\textit{Sprinting}             & 6,5 $\pm$ 1,5  & 5,6 $\pm$ 1,8*  & 2,2 $\pm$ 0,7  & 1,8 $\pm$ 0,8 \\
\hline
\end{tabular}
\end{center}

\begin{itemize}
  \item \textbf{tutti} i parametri di alta intensità sono maggiori negli elite; quelli di bassa
    intensità mostrano differenze piccole e a vantaggio del livello più basso. La prestazione
    d'alto livello si identifica quindi come una prestazione \textbf{a maggiore intensità};
  \item nei \textbf{portieri} la differenza marcata è lo \textbf{stazionamento}: gli elite stanno
    fermi mediamente l'8\% del tempo di gara, gli amatori sensibilmente di più. Entrambi passano
    oltre metà della gara camminando.
\end{itemize}

\rubrica{Risultati --- distanza percorsa nei range di velocità (m)}

\begin{center}
\footnotesize
\begin{tabular}{|l|c|c|c|c|}
\hline
 & \multicolumn{2}{c|}{\textbf{Giocatori di movimento}} & \multicolumn{2}{c|}{\textbf{Portieri}} \\
\hline
\textbf{Distanza percorsa} & \textbf{Elite} & \textbf{Amatori} & \textbf{Elite} & \textbf{Amatori} \\
\hline
\textit{Walking}               & 514 $\pm$ 112   & 551 $\pm$ 127*   & 993 $\pm$ 143  & 960 $\pm$ 126 \\
\textit{Jogging}               & 1302 $\pm$ 671  & 1220 $\pm$ 664*  & 352 $\pm$ 152  & 280 $\pm$ 184 \\
\textit{Low-speed running}     & 1165 $\pm$ 526  & 1019 $\pm$ 573*  & 265 $\pm$ 189  & 189 $\pm$ 127 \\
\textit{Moderate-speed running}& 1050 $\pm$ 355  & 896 $\pm$ 381*   & 196 $\pm$ 130  & 159 $\pm$ 107 \\
\textit{High-speed running}    & 636 $\pm$ 248   & 534 $\pm$ 276*   & 127 $\pm$ 85   & 95 $\pm$ 41 \\
\textit{Sprinting}             & 422 $\pm$ 186   & 308 $\pm$ 203*   & 110 $\pm$ 57   & 87 $\pm$ 46 \\
\hline
\textbf{Distanza totale}       & \textbf{5087 $\pm$ 1104} & \textbf{4528 $\pm$ 1248*} & 2043 $\pm$ 702 & 1770 $\pm$ 854 \\
\hline
\end{tabular}
\end{center}

Gli elite percorrono in termini assoluti circa \textbf{500\,m in più}, e lo scarto si concentra
sulle \textbf{fasce ad alta intensità}: lo \textit{sprinting} passa da 422 a 308\,m.

\rubrica{Risultati --- parametri fisiologici}

\begin{center}
\footnotesize
\begin{tabular}{|l|c|c|c|c|}
\hline
 & \multicolumn{2}{c|}{\textbf{Giocatori di movimento}} & \multicolumn{2}{c|}{\textbf{Portieri}} \\
\hline
\textbf{Parametro} & \textbf{Elite} & \textbf{Amatori} & \textbf{Elite} & \textbf{Amatori} \\
\hline
HR (bpm)                   & 175 $\pm$ 12   & 170 $\pm$ 10*   & 147 $\pm$ 7    & 145 $\pm$ 11 \\
$HR_{max}$ (\%)            & 89,8 $\pm$ 5,8 & 86,2 $\pm$ 6,7* & 73,7 $\pm$ 5,1 & 72,2 $\pm$ 8,6 \\
$VO_2$ (ml/kg/min)         & 43,7 $\pm$ 5,8 & 38,7 $\pm$ 7,9* & 31,5 $\pm$ 4,7 & 29,7 $\pm$ 5,9 \\
$VO_{2}max$ (\%)           & 77,9 $\pm$ 9,0 & 73,1 $\pm$ 6,2* & 63,2 $\pm$ 8,9 & 61,8 $\pm$ 11,7 \\
Dispendio energetico (kcal)& 595 $\pm$ 50   & 543 $\pm$ 67*   & 422 $\pm$ 80   & 415 $\pm$ 65 \\
Lattato ematico (mmol/l)   & 5,5 $\pm$ 1,4  & 5,1 $\pm$ 1,5*  & 4,2 $\pm$ 1,3  & 4,0 $\pm$ 1,9 \\
\hline
\end{tabular}
\end{center}

I valori degli elite \textbf{confermano quelli di Castagna e D'Ottavio}: FC al 90\% circa della
massima, $VO_2$ al 78\% contro il 76\%, lattato 5,5 contro 5,3\,mmol/l. Si può quindi affermare con
ragionevole certezza che la prestazione nel calcio a 5 di buon livello si attesta \textbf{intorno al
90\% della frequenza cardiaca massima}.

\rubrica{Il tempo speso ad alta intensità}
Percentuale del \textbf{tempo di gara} trascorsa sopra l'85\% della $HR_{max}$: elite
\textbf{81,4 $\pm$ 16,3\%}, amatori \textbf{73,5 $\pm$ 21,4\%}.

\rubrica{Perché questa intensità è sostenibile}
Valori così alti sono possibili grazie a una regola specifica del futsal (cap.~\ref{cap:introfutsal}):
si può \textbf{rientrare in campo dopo essere stati sostituiti}. Nell'alta qualificazione i
giocatori restano in campo \textbf{3--5 minuti effettivi}, poi vengono sostituiti, recuperano e
rientrano pronti a sostenere di nuovo quel livello di intensità.

\rubrica{Take home message}
\begin{itemize}
  \item il \textbf{futsal d'elite richiede un alto impegno fisiologico};
  \item differenze fra elite e non elite nei parametri fisiologici: \textbf{HR}, \textbf{$VO_2$},
    \textbf{lattato};
  \item differenze fra elite e non elite nella \textbf{distanza percorsa ad alta intensità}.
\end{itemize}

% ---------------------------------------------------------------------------
\section{Partita ufficiale e amichevole a confronto}

\subsection{Vieira e coll., 2016}

\textit{Preliminary results on organization on the court, physical and technical performance of
Brazilian professional futsal players: comparison between friendly pre-season and official match},
Motriz, Rio Claro, 22(2), 79--91.

\rubrica{Perché interessa}
Negli sport di squadra la valutazione funzionale --- specie del metabolismo aerobico --- si svolge
generalmente quando il campionato è fermo: \textbf{prima dell'inizio}, durante la preparazione
precampionato, nell'eventuale \textbf{pausa di metà anno} e a \textbf{fine stagione}. La domanda è
quindi se una partita amichevole di precampionato possa essere usata per valutare il giocatore.

\rubrica{Metodo e protocollo}
Video tracking; \textbf{1 partita ufficiale} (OM) e \textbf{1 amichevole} (FM), $n = 10$ giocatori
--- gli \textbf{stessi} giocatori nelle due condizioni, non due livelli di qualificazione diversi.
Il numero di rilevazioni è \textbf{estremamente basso}: lo studio serve a farsi un'idea, non a
concludere. Range di velocità di Castagna e D'Ottavio; valori in \textbf{mediana (IQR)}.

\rubrica{Risultati --- performance fisica}

\begin{center}
\footnotesize
\begin{tabular}{|l|c|c|}
\hline
\textbf{Variabile (partita intera)} & \textbf{Amichevole (FM)} & \textbf{Ufficiale (OM)} \\
\hline
TD --- distanza totale (m)          & 2654,48 (1854,13) & 3191,13 (1455,50) \\
$t$ --- tempo in campo (min)        & 29,82 (21,50)     & 29,66 (20,56) \\
NS/$t$ --- sprint per tempo (a.u.)  & 0,73 (0,27)       & 0,92 (0,64) \\
$V_{MAX}$ (km/h)                    & 27,67 (5,50)      & 28,16 (3,36) \\
$D_{MIN}$ (m/min)                   & 94,79 (10,14)     & 103,53 (19,68)\textsuperscript{\#a} \\
\hline
\end{tabular}
\end{center}
\footnotesize \textsuperscript{\#}\,valori significativamente diversi fra partita ufficiale e
amichevole; \textsuperscript{a}\,\textit{large effect size} rispetto all'amichevole. NS: numero di
sprint; a.u.: unità arbitrarie.
\normalsize

\begin{itemize}
  \item la \textbf{distanza totale} è nettamente superiore in partita ufficiale, con circa
    \textbf{500\,m di differenza}. I 3191\,m sono peraltro \textbf{equiparabili} ai circa 3100 di
    Bueno;
  \item il \textbf{numero di sprint} rapportato al tempo è nettamente superiore in gara ufficiale;
  \item la \textbf{velocità massima} tocca i 28\,km/h in gara ufficiale, valore molto elevato, e si
    avvicina in amichevole. In \textbf{entrambe} le condizioni la massima del secondo tempo è
    inferiore a quella del primo --- al contrario della distanza percorsa, che in termini assoluti
    è maggiore nel secondo tempo;
  \item la \textbf{distanza al minuto} si attesta a $\approx 105$\,m/min in gara ufficiale, contro
    i 120 di Castagna e gli oltre 130 dell'\textit{in play} di Bueno.
\end{itemize}

\rubrica{Risultati --- ripartizione nelle bande di velocità (\% della distanza, partita intera)}

\begin{center}
\footnotesize
\begin{tabular}{|l|c|c|}
\hline
\textbf{Banda} & \textbf{Amichevole (FM)} & \textbf{Ufficiale (OM)} \\
\hline
$V_1$ & 35,67 (6,10) & 30,38 (8,24)\textsuperscript{\#a} \\
$V_2$ & 39,49 (3,07) & 41,27 (3,36) \\
$V_3$ & 12,99 (1,50) & 16,41 (2,21)\textsuperscript{\#a} \\
$V_4$ & 6,13 (1,48)  & 7,51 (1,72) \\
$V_5$ & 6,15 (3,55)  & 5,27 (5,46) \\
\hline
\end{tabular}
\end{center}

\begin{itemize}
  \item in amichevole i giocatori trascorrono il \textbf{35\%} della distanza a bassissima velocità
    ($V_1$), percentuale che scende drasticamente in gara ufficiale: è una delle due differenze
    significative, insieme alla crescita della media intensità ($V_3$);
  \item il maggiore stazionamento del secondo tempo, statisticamente significativo, si può
    attribuire all'intervento della \textbf{fatica muscolare} e interpretare come necessità di un
    \textbf{maggior tempo di recupero};
  \item lo \textbf{sprint} ($V_5$) fa eccezione: il trend è \textbf{leggermente superiore in
    amichevole}. È come se in amichevole il giocatore avesse comunque la \textbf{disponibilità
    bioenergetica} per gesti ad alta intensità, nonostante la scarsa preparazione fisico-atletica o,
    al contrario, una risposta adattativa non ancora adeguata ai carichi elevati del precampionato.
\end{itemize}

\rubrica{Risultati --- performance tecnica}
Frequenza di esecuzione dei gesti tecnici, normalizzata sul tempo in campo, in media $\pm$
deviazione standard.

\begin{center}
\footnotesize
\begin{tabular}{|l|c|c|c|c|}
\hline
 & \multicolumn{2}{c|}{\textbf{1\textsuperscript{o} tempo}} & \multicolumn{2}{c|}{\textbf{Partita intera}} \\
\hline
\textbf{Gesto (al minuto)} & \textbf{FM} & \textbf{OM} & \textbf{FM} & \textbf{OM} \\
\hline
Passaggi          & 0,9 $\pm$ 0,29  & 1,9 $\pm$ 0,38\textsuperscript{*\#ab} & 0,9 $\pm$ 0,29  & 1,45 $\pm$ 0,3\textsuperscript{\#b} \\
Possessi          & 0,97 $\pm$ 0,25 & 1,9 $\pm$ 0,33\textsuperscript{*\#ab} & 1 $\pm$ 0,15    & 1,48 $\pm$ 0,32\textsuperscript{\#b} \\
Tocchi di palla   & 2,29 $\pm$ 0,9  & 5 $\pm$ 1,43\textsuperscript{*\#ab}   & 2,31 $\pm$ 0,63 & 3,85 $\pm$ 0,94\textsuperscript{\#b} \\
Duelli            & 0,45 $\pm$ 0,17\textsuperscript{\#c} & 0,24 $\pm$ 0,12 & 0,34 $\pm$ 0,14 & 0,29 $\pm$ 0,11 \\
Tiri              & 0,07 $\pm$ 0,06 & 0,1 $\pm$ 0,09  & 0,08 $\pm$ 0,06 & 0,11 $\pm$ 0,06 \\
\hline
\end{tabular}
\end{center}
\footnotesize \textsuperscript{*}\,differenza significativa fra primo e secondo tempo;
\textsuperscript{\#}\,differenza significativa fra ufficiale e amichevole;
\textsuperscript{a}\,\textit{large effect size} rispetto al secondo tempo;
\textsuperscript{b}\,rispetto all'amichevole; \textsuperscript{c}\,rispetto alla partita ufficiale.
\normalsize

Le tre voci su cui il confronto è netto sono \textbf{passaggi}, \textbf{possessi} e \textbf{tocchi
di palla} al minuto: nella partita ufficiale il primo tempo fa registrare valori \textbf{circa
doppi} rispetto all'amichevole, e il vantaggio si dimezza nel secondo tempo. Anche il parametro
tecnico, quindi, subisce un decremento per effetto della \textbf{fatica}.

\rubrica{Risultati --- l'esecuzione riuscita}
Gli autori hanno valutato non solo il numero di eventi tecnici ma anche la loro \textbf{qualità}:
\begin{itemize}
  \item i \textbf{passaggi corretti} al minuto passano da 1,75 $\pm$ 0,43 nel primo tempo della
    partita ufficiale a 0,95 $\pm$ 0,39 nel secondo: un \textbf{forte decremento}. In amichevole il
    trend è opposto, anche se di intensità minore (0,74 e 0,85);
  \item i \textbf{duelli vinti} al minuto si comportano diversamente. Questa qualità --- dribbling,
    finta, marcamento e smarcamento --- dà valori \textbf{simili} fra le due tipologie di partita
    (0,19 su entrambe, sulla partita intera), e il trend fra i due tempi è \textbf{esattamente
    opposto}: in amichevole se ne vincono di più nel primo tempo (0,29 contro 0,14), in partita
    ufficiale di più nel secondo (0,12 contro 0,25).
\end{itemize}

\rubrica{Take home message}
\begin{itemize}
  \item \textbf{differenze significative} fra amichevole e partita ufficiale per l'\textbf{alta
    intensità} e per i \textbf{parametri tecnici};
  \item la \textbf{performance tecnica è influenzata dalla fatica muscolare} --- ma non tutta: la
    componente dei duelli non segue questo andamento;
  \item le \textbf{amichevoli non vanno usate} per valutare il giocatore, né dal punto di vista
    fisico-atletico né da quello tecnico, né per giudicarne lo stato di forma, né come parametro di
    riferimento per l'allenamento. In preparazione è preferibile svolgere \textbf{test di
    valutazione codificati}.
\end{itemize}

% ---------------------------------------------------------------------------
\section{Il futsal femminile}

\subsection{Beato e coll., 2017}

\textit{Evaluation of the external and internal workload in female futsal players}, Biology of
Sport, 34(3), 227--231.

Il futsal femminile ha \textbf{pochissimi studi} nella letteratura scientifica internazionale:
questo è recente e, pur con i limiti che seguono, dà un'idea della prestazione.

\rubrica{Perché qui il GPS si può usare}
È l'unico studio della serie a non impiegare il video tracking: la partita è stata giocata in
condizione \textbf{outdoor}, su un campo in sintetico di dimensioni regolamentari, e questo ha
permesso l'uso del \textbf{GPS a 10\,Hz}.

\rubrica{Campione e limiti}
$n = 16$ giocatrici: età 27 $\pm$ 5 anni, statura 1,65 $\pm$ 0,09\,m, massa corporea
56,9 $\pm$ 7,7\,kg, massa grassa 21,5 $\pm$ 2,9\%. Due limiti da tenere presenti nel confronto con
gli studi maschili:
\begin{enumerate}
  \item il tempo \textbf{non è stato fermato} quando la palla usciva: non è tempo effettivo. Si
    sono giocati due tempi da 20 minuti;
  \item le giocatrici sono rimaste in campo \textbf{per tutto il tempo}, senza rotazioni: il loro
    tempo di gioco è quindi molto superiore a quello dei giocatori maschi degli altri studi.
\end{enumerate}

\rubrica{Risultati}
Valori in media $\pm$ deviazione standard; le differenze in valore assoluto e in percentuale.

\begin{center}
\footnotesize
\begin{tabular}{|l|c|c|c|c|}
\hline
\textbf{Attività locomotoria} & \textbf{1\textsuperscript{o} tempo} & \textbf{2\textsuperscript{o} tempo} & \textbf{Differenza} & \textbf{Partita} \\
\hline
Tempo                        & 20\,min\,25\,s & 20\,min\,26\,s & ---          & 40\,min\,51\,s \\
TD --- distanza totale (m)   & 1424 $\pm$ 114* & 1313 $\pm$ 113 & $-111$ (7,8\%)  & 2737 $\pm$ 207 \\
RV --- velocità relativa (m/min) & 70 $\pm$ 6* & 64 $\pm$ 6     & $-4$ (8,6\%)    & 67 $\pm$ 5 \\
HSR --- \textit{high speed running} (m) & 28 $\pm$ 16* & 22 $\pm$ 19 & $-6$ (21,4\%) & 50 $\pm$ 33 \\
HMD --- \textit{high metabolic distance} (m) & 80 $\pm$ 29* & 69 $\pm$ 26 & $-11$ (13,7\%) & 150 $\pm$ 53 \\
Numero di accelerazioni      & 16 $\pm$ 4      & 16 $\pm$ 4     & 0 (0\%)         & 31 $\pm$ 6 \\
Numero di decelerazioni      & 21 $\pm$ 6      & 19 $\pm$ 7     & $-2$ (9,5\%)    & 40 $\pm$ 12 \\
DBL --- \textit{dynamic body load} (AU) & 49 $\pm$ 26 & 51 $\pm$ 28 & $+2$ (4,1\%) & 101 $\pm$ 55 \\
MP --- potenza metabolica media (W/kg) & 6,2 $\pm$ 0,7\textsuperscript{\#} & 5,9 $\pm$ 0,7 & $-0,3$ (4,8\%) & 6,1 $\pm$ 0,6 \\
HR media (\% della massima)  & 85 $\pm$ 3      & 81 $\pm$ 4     & $-4$ (4,7\%)    & 83 $\pm$ 3 \\
\hline
\end{tabular}
\end{center}
\footnotesize \textsuperscript{*}\,$p < 0{,}05$ fra primo e secondo tempo;
\textsuperscript{\#}\,$p = 0{,}059$.
\normalsize

\begin{itemize}
  \item la \textbf{distanza totale} di 2737\,m è vicina a quella dei colleghi uomini --- ma le
    giocatrici sono rimaste in campo molto più a lungo;
  \item la \textbf{distanza al minuto} lo rivela: 70\,m/min nel primo tempo, 64 nel secondo, 67
    sulla partita, contro i 105, 120 e oltre 130\,m/min dei tre studi maschili;
  \item la \textbf{frequenza cardiaca} si attesta all'\textbf{83\%} della massima, contro il 90\%
    circa degli uomini;
  \item il calo fra primo e secondo tempo interessa tutte le variabili di intensità ed è
    percentualmente \textbf{più marcato sulle fasce più intense} ($-21{,}4\%$ per l'HSR,
    $-13{,}7\%$ per l'HMD, contro $-7{,}8\%$ per la distanza totale): è la testimonianza della
    \textbf{fatica muscolare} sopravvenuta, coerente con il fatto che non siano mai uscite dal
    campo. Fanno eccezione il \textbf{numero di accelerazioni}, identico nei due tempi, e il
    \textbf{dynamic body load}, che aumenta.
\end{itemize}

\rubrica{Take home message}
La prestazione femminile si avvicina a quella maschile come \textbf{volume totale}, ma \textbf{non
per l'alta intensità}: da questo punto di vista non è equiparabile. Allenando una squadra femminile
occorre tenerne conto nella programmazione e nella scelta dei mezzi di allenamento.

% ---------------------------------------------------------------------------
\section{Il quadro d'insieme}

\rubrica{La distanza al minuto}
È il parametro più utile per l'allenamento, e cambia molto a seconda di che cosa si misura:
\begin{itemize}
  \item Bueno, \textit{in play}: oltre \textbf{130}\,m/min --- il valore reale del gioco effettivo;
  \item Castagna: \textbf{121}\,m/min;
  \item Vieira, partita ufficiale: $\approx$ \textbf{105}\,m/min;
  \item Bueno, partita intera: $\approx$ \textbf{90--98}\,m/min --- il valore che si ottiene
    ignorando la distinzione \textit{in play};
  \item Beato, femminile: \textbf{67}\,m/min.
\end{itemize}

\rubrica{I valori fisiologici convergono}
Due studi indipendenti danno valori sovrapponibili per il giocatore di buon livello: FC intorno al
\textbf{90\%} della massima, $VO_2$ al \textbf{76--78\%} del massimo, lattato \textbf{5,3--5,5}\,mmol/l.

\rubrica{Che cosa resta}
Il calcio a 5 è un \textbf{regime intermittente ad alta intensità}, in cui il lattato prodotto viene
riutilizzato nei recuperi, la \textbf{RSA} è la capacità specifica determinante e il secondo tempo
mostra costantemente i segni della \textbf{fatica} --- sul piano cinematico e anche su quello
tecnico. Conoscere questi valori è la premessa per scegliere i test di valutazione e per
\textbf{parametrare i loro risultati} all'impegno reale che il giocatore sostiene in partita.
